% !PRINTED 153
\weberpage{153}
%<<CONT
$j(\omega)$ bezeichnen und schlechtweg die \gesperrt{Invariante} nennen und so definieren:
\weq{15}{j(\omega) = 2^8\frac{(1-\varkappa^2\varkappa'^2)^3}{\varkappa^4\varkappa'^4},}
woraus wir erhalten:
\weq{16}{\begin{aligned}
\frac{4.27g_2^3}{G} &= j(\omega), \\
\frac{4.27.27g_3^2}{G} &= j(\omega)-27.64 = \frac{64.(2+\varkappa^2\varkappa'^2)^2(\varkappa'^2-\varkappa^2)^2}{\varkappa^4\varkappa'^4}.
\end{aligned}}
Als Funktion von $\varkappa^2$ betrachtet, hat die Funktion $j(\omega)$ die Eigenschaft,
ungeändert zu bleiben, wenn $\varkappa^2$ durch einen der sechs Werte
\[ \varkappa^2, \quad \varkappa'^2, \quad \frac{1}{\varkappa^2}, \quad \frac{1}{\varkappa'^2}, \quad -\frac{\varkappa^2}{\varkappa'^2}, \quad -\frac{\varkappa'^2}{\varkappa^2} \]
ersetzt wird.

Wenn wir nach (7) den Differentialquotienten der Funktion $\wp(u)$ bilden, so erhalten wir
\weq{17}{\wp'(u) = -\frac{16K^3}{\omega_1^3}\frac{\mathrm{cn}\,v\,\mathrm{dn}\,v}{\mathrm{sn}^3v} = -2\frac{\sigma_{00}(u)\sigma_{10}(u)\sigma_{01}(u)}{\sigma(u)^3}}
und daraus nach (3):
\weq{18}{\begin{aligned}
\wp'(u)^2 &= 4[\wp(u)-e_1][\wp(u)-e_2][\wp(u)-e_3] \\
&= 4\wp(u)^3-g_2\wp(u)-g_3,
\end{aligned}}
oder endlich
\weq{19}{du = \frac{d\wp}{\sqrt{4\wp^3-g_2\wp-g_3}},}
woraus man ersieht, daß die Funktion $\wp(u)$ in derselben Beziehung zur
\gesperrt{Weierstrassschen} Normalform des elliptischen Differentials steht, wie die
Funktion $\mathrm{sn}\,v$ zu der \gesperrt{Legendreschen}.

\begin{center}
\textbf{§\,47. \gesperrt{Die elliptischen Transzendenten zweiter Gattung.}}
\end{center}

\gesperrt{Jacobi} hat als Transzendente zweiter Gattung die Funktion
\weq{1}{Z(v) = \frac{\Theta'(v)}{\Theta(v)} = \frac{d\log\Theta(v)}{dv}}
eingeführt. Die Beziehung dieser Funktion zu den elliptischen Integralen zweiter Gattung
ergibt sich aus der Formel (16) des §\,23, wobei gleich bemerkt sei, daß ganz ähnliche
Betrachtungen an die dortigen Formeln (15), (17), (18) anzuknüpfen wären, die aber nicht
zu wesentlich neuen Resultaten führen.
